\documentclass{article}
\usepackage[utf8]{inputenc}
\usepackage[russian]{babel}
\usepackage{amsmath}
\usepackage{amssymb}
\usepackage{mathtools}
\usepackage[a4paper, left=1.5cm, right=1.5cm, top=1.5cm, bottom=1.5cm]{geometry}\setlength{\parindent}{0pt}

\title{Билет 17. Вычисление \( \int_{0}^{+\infty} \frac{\sin x}{x} \, dx \).}
\author{}
\date{}

\begin{document}

\maketitle

\[
\int_{0}^{+\infty} \frac{\sin x}{x} \, dx \quad - \text{ интеграл Дирихле}
\]

\[
\downarrow \quad \int \frac{\sin x}{x} \, dx \text{ — не берущийся.}
\]

Несобственный интеграл 1 рода, поскольку при \( x = 0 \) функция не имеет особенности.
\[
\lim_{x \to 0} \frac{\sin x}{x} = 1,
\]
следовательно, это не несобственный интеграл 2 рода.

Свойство интегралов Римана: если функция переопределена в конечном числе точек,
это не мешает сходимости интеграла.

\[
f(x) =
\begin{cases}
\frac{\sin x}{x}, & x > 0 \\
1, & x = 0
\end{cases}
\quad \text{Доопределение на } x=0 \text{ не влияет на значение интеграла.}
\]

\[
[0, +\infty) \text{ — функция непрерывна.}
\]

\text{Применим принцип Дирихле для сходимости интеграла:}

\[
h(x) = \sin x \quad g(x) = \frac{1}{x}
\]

\[
F(x) = \cos x \quad \text{ограничена.}
\]

\text{Все условия выполнены. Интеграл сходится. Теперь нужно найти его значение.}

\[
\int_{0}^{+\infty} \frac{\sin(\lambda x)}{x} \, dx \quad \lambda \geq 0
\]

\text{Опять доопределим функцию:}

\[
f(x) =
\begin{cases}
\frac{\sin(\lambda x)}{x}, & x > 0 \\
\lambda, & x = 0
\end{cases}
\quad \lim_{x \to 0} \frac{\sin(\lambda x)}{x} = \lambda
\]

Функция непрерывна. Интеграл сходится от параметра,
но пока не можем работать с интегралом, введём ещё один параметр.

\[
\int_{0}^{\infty} e^{-\gamma x} \cdot \frac{\sin(\lambda x)}{x} \, dx, \quad \gamma > 0, \quad x \geq 0
\]

Признак Абеля гарантирует сходимость.
Теперь дифференцируем по \( \lambda \) (теорема 4).
После этого перейдём к пределу (теорема 1) при \( \gamma \to 0 \).

\[
\varphi(x, \lambda) = e^{-\gamma x} \frac{\sin(\lambda x)}{x}
\]

\[
\varphi(x, 0) = 0, \quad \varphi(0, \lambda) = \lambda \quad \checkmark
\]

1) \( \varphi(x, \bar{\lambda}) \) — непрерывна на \( [0, +\infty) \), выполняются условия для дифференцирования под знаком интеграла.

2) \( \frac{\partial \varphi(x, \lambda)}{\partial \lambda} = e^{-\gamma x} \cos(\lambda x) \)

3) Применяем принцип сравнения.

4) \( \int_{0}^{+\infty} e^{-\gamma x} \cos(\lambda x) \, dx \quad |e^{-\gamma x} \cos(\lambda x)| \leq e^{-\gamma x} \)

\[
\int_{0}^{+\infty} e^{-\gamma x} \, dx \quad \text{сходится.}
\]

\text{Таким образом, все 4 условия для дифференцирования под знаком интеграла выполнены.}

\text{Теперь можно дифференцировать.}

\[
G'(\lambda) = \int_{0}^{+\infty} e^{-\gamma x} \cos(\lambda x) \, dx
\]

\text{Теперь применяем интегрирование по частям. Рассмотрим:}

\[
I = \int_{0}^{\infty} e^{-\gamma x} \cos(\lambda x) \, dx
\]

Решение по частям:

\[
I = \left[ \frac{e^{-\gamma x} \sin(\lambda x)}{\lambda} \right]_{0}^{\infty} - \int_{0}^{\infty} \frac{e^{-\gamma x} \cos(\lambda x)}{\lambda} \, dx
\]

На границе \(x = \infty\) вторая часть стремится к нулю, поскольку \(e^{-\gamma x} \to 0\) при \(x \to +\infty\), а \( \sin(\lambda x) \) и \( \cos(\lambda x) \) ограничены. На границе \(x = 0\) мы получаем:

\[
\frac{e^{-\gamma x} \sin(\lambda x)}{\lambda} \Big|_0^\infty = 0
\]

Таким образом, остается:

\[
I = \int_{0}^{\infty} \frac{\gamma}{\gamma^2 + \lambda^2} \, dx = \frac{1}{\gamma^2 + \lambda^2}
\]

\text{Теперь используем формулу Ньютона-Лейбница:}

\[
G(\lambda) - G(0) = \int_{0}^{\lambda} \frac{\gamma}{\gamma^2 + t^2} \, dt = \arctg \frac{t}{\gamma} \Bigg|_{0}^{\lambda} =
\]

\[
= \arctg \frac{\lambda}{\gamma} \quad * G(0) = 0
\]

\[
! \text{ Так как } \lambda \geq 0, \gamma > 0, \quad \text{то } \lambda \geq 0 \text{ можно подставить.}
\]

\[
\text{А } \gamma = 0 \quad \text{не можем / ну и делать нельзя.}
\]

Рассмотрим:

\[
H(\gamma) = \int_{0}^{\infty} e^{-\gamma x} \frac{\sin(\lambda x)}{x} \, dx = \arctg \frac{\lambda}{\gamma}
\]

\[
\gamma > 0, \lambda > 0; \quad \lambda \text{ зафиксирован, рассматриваем зависимость от } \gamma.
\]

Теперь переходим к пределу:

\[
\lim_{\gamma \to 0} H(\gamma) = \frac{\pi}{2}
\]

Переходя к пределу под знаком интеграла, получаем:

\[
\text{Условия теоремы:}
\]
1) \( \forall A > 0: \frac{e^{-\gamma x} \sin(\lambda x)}{x} \xrightarrow[\gamma \to +0]{} \frac{\sin(\lambda x)}{x} \) равномерно на промежутке \( [0, A] \).

\[
\left| e^{-\gamma x} \frac{\sin(\lambda x)}{x} - \frac{\sin(\lambda x)}{x} \right| = \left| \frac{\sin(\lambda x)}{x} \right| \cdot \left| e^{-\gamma x} - 1 \right| \leq \lambda A \gamma \xrightarrow[\gamma \to +0]{} 0
\]

\[
e^{-z} - 1 = q(z) - q(0) =
\]

Если \( q = e^{-z} \), то \( q \) дифференцируема,
т.е. применима формула конечных приращений Лагранжа.

\[
\ominus 0 < \Theta < 1
\]

\[
= z \cdot q'(\Theta z) = -z \cdot e^{-\Theta z} 
\]

\[
\quad \left| e^{-\gamma x} - 1 \right| = |\gamma x| e^{-\Theta \gamma x} \leq |\gamma x| \leq |\gamma| A
\]

\[
x \in [0, A]
\]

\text{Но еще было условие, что интеграл равномерно сходится на множестве \( Y \).}

Пример по критерию Коши: нужно показать, что
\[
\forall \varepsilon > 0 \ \exists B :
\forall A', A'' \geq B \ \forall \gamma > 0
\]

\[
\left| \int_{A'}^{A''} e^{-\gamma x} \frac{\sin(\lambda x)}{x} \, dx \right| < \varepsilon
\]

\[
\lim_{\gamma \to 0} H(\gamma) = \frac{\pi}{2} = \int_{0}^{\infty} \frac{\sin(\lambda x)}{x} \, dx, \quad \lambda > 0
\]

\[
\boxed{\text{Ответ: } \frac{\pi}{2}.}
\]

\end{document}
